\subsection{4.3 组内智能体评分}\label{groupwise-agentic-grading}

对于代码智能体任务，二值测试奖励提供了可扩展的正确性信号，但无法区分通过解之间的实现质量或解题行为。我们在这些任务的不同子集上使用两种互补的方法。组内奖励合成（GRS）用于高通过率任务的子集，它比较多个离线 rollout，构造任务专属的评分细则，这些细则在训练期间被复用来为单个 rollout 打分，并将其质量分与测试奖励结合。对于其余全部代码智能体任务，我们主要依赖组内优势重分配（GAR）。其在线智能体评分器联合考察每个混合结果组内的成功与失败轨迹，对通过解排序，并将正向优势重新分配给质量更高的通过轨迹。图 7 概括了这两条工作流。




\begin{figure}[htbp]
\centering
\includegraphics[width=\linewidth]{images_hi/fig07.png}
\caption{用于代码智能体 RL 的组内智能体评分。(a) 组内奖励合成将测试结果与预计算的任务专属评分细则给出的逐 rollout 分数相结合。(b) 组内优势重分配在线比较轨迹，将确认作弊的奖励重置为零，并重分配序列级优势。柱状图示意正向优势从质量较低的通过解向质量较高的通过解的转移。}
\label{fig:7}
\end{figure}

\subsection{4.3.1 组内奖励合成（GRS，离线评分细则）}\label{groupwise-reward-synthesis-grs-offline-rubrics}

对于每个选定的任务，我们收集多个离线 rollout，并要求智能体将它们与任务说明和代码仓库一起研读。比较这些尝试可以揭示不同的解法、反复出现的错误，以及可能分散在若干轨迹中的有用行为。智能体将这一分析转化为两组标准：解法评分细则，用于评估最终的实现；以及行为评分细则，用于评估智能体如何着手并验证其工作。

解法评分细则描述良好实现中与任务相关的性质，包括是否满足需求、是否恰当地处理边界情况，以及与周围代码库一致的改动。行为评分细则描述支持可靠解题的可观察做法，例如收集相关证据、检查代码改动的效果。我们寻求的标准要能区分质量上的实质差异，同时允许不同的有效解题方式。

采样得到的 rollout 为评分细则的构造提供依据，但标准本身植根于任务。一个有用的行为可能只出现在一次尝试中，而任务所支持的一项需求可能在任何已观察到的解法中都不存在。因此智能体能够识别出超出采样 rollout 所展示的改进点。同时，某一个成功解法所做的选择不会自动成为其他所有解法的要求。

得到的评分细则被复用来逐一评估后续的训练 rollout。评分器智能体进入每个 rollout 的执行环境，依据任务专属的评分细则对其进行评估，并以生成的代码、执行结果和轨迹作为证据。它给出一个解法分数


\(S ^ { \mathrm { s o l } }\) 以衡量实现质量，并给出一个行为分数 \(S ^ { \mathrm { b e h } }\) 以衡量智能体如何着手并验证其工作。

对于轨迹 $i$，记 \(R _ { i } ^ { \mathrm { t e s t } }\) 为原始的二值测试奖励。我们通过将该测试奖励与两个评分细则分数相乘来合成最终训练奖励：

\[
R _ { i } = R _ { i } ^ { \mathrm { t e s t } } \cdot S _ { i } ^ { \mathrm { s o l } } \cdot S _ { i } ^ { \mathrm { b e h } } .\tag{2}
\]

这种相乘形式使评分细则的监督与测试结果绑定。失败的轨迹保持零奖励，而通过的轨迹则按实现质量和解题行为进一步区分。即便一个组内的每个 rollout 都通过了测试，两个评分细则分数之积的差异仍可提供学习信号。离线任务分析由此成为可复用的监督来源，使训练能够捕捉二值测试奖励未能表达的质量差异。

\subsection{4.3.2 组内优势重分配（GAR，在线评分）}\label{groupwise-advantage-redistribution-gar-online-grading}

我们对其余代码智能体任务应用在线组内优势重分配。对于每个混合结果的 rollout 组，我们将所有轨迹放入一个共享工作区，其中包含任务说明、代码仓库、提交的补丁和测试输出。一个经 SFT 训练的智能体评分器联合考察组内的所有轨迹，对比成功与失败的尝试，并从五个维度比较通过的补丁：解法思路的合适性，实现该思路时在没有遗漏或不必要回退上的精确性，相对于必要改动的最小性，是否避免任务之外的非预期影响，以及与代码库约定一致的工程质量。这些评估指导组内排序，当差异不具决定性时允许并列。评分器可以检查仓库代码并运行针对性测试，以更好地理解任务并验证候选解法。当证据确认某条轨迹依赖外部答案或泄露答案时，我们将该轨迹的有效奖励重置为零，并在重算组内统计量之前将其视为失败。

我们使用剩余的质量排序来重分配序列级优势。对于轨迹 \(i ,\)，记 \(R _ { i }\) 为作弊修正后的有效二值奖励，$\bar{R}$ 为其组均值，\(A _ { i } = R _ { i } - \bar { R }\) 为其序列级优势，\(\mathcal { P } \ = \ \{ i \ : \ R _ { i } \ = \ 1 \}\) 。对于非空的 \(\mathcal { P } _ { z }\) 质量因子 \(f _ { i } ~ \in ~ ( 0 , 1 ]\) 首先降低质量较低的通过解的权重，随后一个公共因子将被移除的正向优势质量在通过轨迹之间重分配。未加封顶的更新为

\[
\lambda = \frac { \sum _ { j \in \mathcal { P } } A _ { j } } { \sum _ { j \in \mathcal { P } } f _ { j } A _ { j } } , \qquad A _ { i } ^ { \prime } = \left\{ \begin{array} { l l } { \lambda f _ { i } A _ { i } , } & { i \in \mathcal { P } , } \\ { A _ { i } , } & { i \notin \mathcal { P } . } \end{array} \right.\tag{3}
\]

这种未加封顶的更新保持
\(\begin{array} { r } { \sum _ { i \in \mathcal { P } } A _ { i } ^ { \prime } \ = \ \sum _ { i \in \mathcal { P } } A _ { i } } \end{array}\)
以及由质量诱导的相对权重不变，同时不改变失败的轨迹。实际上，它把正向优势质量从质量较低的通过轨迹重分配给质量较高的通过轨迹。仅降低正向优势的权重会让负向优势保持不变；重归一化则恢复二者的平衡，作为防止熵过度增长的保障。实践中，我们对公共缩放因子加以封顶，以防止正向优势被过度放大。随后我们从通过和失败两类轨迹的优势中减去组均值，得到组均值为零的最终序列优势 \(A _ { i } ^ { \mathrm { n e w } }\)。我们将该序列级优势广播到轨迹中的所有回复 token。评分异步运行，不可用的评分器输出回退到原始优势。

为评估在线组内优势重分配对训练动态的影响，我们比较 MiMo-V2.6-Flash 在有无该方法下的纯代码 RL 运行，训练批大小为 128，损失按 token 均值聚合（图 8）。没有在线评分时，轮数和总 token 长度迅速增长，导致更多轨迹触及长度上限，使通过率的提升难以持续。有在线评分时，通过率的增益持续到第 52 步，同时轮数大致稳定，token 长度缓慢增长。这些趋势表明，在线评分支撑了在更稳定的训练状态下持续的策略改进。



\begin{figure}[htbp]
\centering
\includegraphics[width=\linewidth]{images_hi/fig08.png}
\caption{在 DeepSWE v1.1 上有无在线组内优势重分配的对比，使用 MiMo-V2.6-Flash 的纯代码 RL（训练批大小 128）。通过率为 avg@n，其中 \(n = 3 ;\) 我们还报告了平均总轮数和总 token 长度。}
\label{fig:8}
\end{figure}


另外的面向维护者的审计发现，在提高测试通过率的压力下，未经在线评分训练的策略越来越多地采用不良行为，例如投机性的兼容分支、大范围的导出、吞掉异常、放松校验，以及针对评测的配置改动。这些变通做法旨在提高通过测试的可能性，但可能超出任务指令的范围、不必要地扩大 API、掩盖失败，并使代码更难维护。相比之下，经在线评分训练的策略倾向于产生更小、更精确的补丁，这些补丁保持在所要求的范围内，也更容易维护。
